\SourceLeaf{364}%
einfach durch \(240:1823\) vorstellen liess; diesen Vortheil hätte er entbehren müssen, wäre er ursprünglich von dem Bruche \(\frac{362}{209}\) ausgegangen.

Am Schlusse der Abhandlung macht Hr. M. noch die Bemerkung, dass auch \textsc{Columella} \textit{de re rustica} V, 2 von einem der genäherten Werthe von \(\BareRoot{3}\) (nemlich von \(\frac{26}{15}\)) Gebrauch gemacht hat, indem er für den Inhalt des gleichseitigen Dreiecks die Summe des dritten und des zehnten Theils des auf seiner Seite beschriebenen Quadrats annimmt.

\pdfbookmark[0]{Anzeigen. Monge: Géométrie descriptive. 1813.}{anzeigen-monge-1813}
\begin{center}
\rule{.85\textwidth}{.3pt}\par
Göttingische gelehrte Anzeigen.\quad 1813 Juli 31.\par
\rule{.85\textwidth}{.3pt}
\end{center}
\hspace*{\parindent}\textit{Géométrie descriptive par} \textsc{Gaspard Monge}, \textit{de l’institut des sciences etc. Nouvelle édition. Avec un supplément par} \textsc{M. Hachette}, \textit{instituteur à l’école impériale polytechnique etc.} Paris, bei \textsc{J. Klostermann} dem jüngern. 162 und 118 Seiten in Quart.

Die Geometrie, deren Gegenstand die Raumverhältnisse sind, zerfällt in zwei grosse Abtheilungen, je nachdem der Raum nur nach zwei Dimensionen betrachtet wird (in der Ebene), oder nach allen drei Dimensionen zugleich. Man begreift leicht, dass der andere Theil seiner Natur nach von einem viel grössern Umfange sein, und eine viel grössere Mannigfaltigkeit von Fragen und Untersuchungen darbieten müsse, als der erste. Wenn daher schon von unserer Elementar-Geometrie die Planimetrie einen grössern Theil ausmacht, als die Stereometrie, so rührt dies nur daher, dass letztere verhältnissmässig viel weniger entwickelt und ausgebildet ist. In der That hat man vorzüglich die Untersuchungen der letztern Art in neuern Zeiten lieber mit Hülfe der Analyse behandelt, und sie so gleichsam der Geometrie entzogen, welche sich nur der unmittelbaren Anschauung bedient. Es ist auch nicht zu läugnen, dass die Vorzüge der analytischen Behandlung vor der geometrischen, ihre Kürze, Einfachheit, ihr gleichförmiger Gang, und besonders ihre Allgemeinheit, sich gewöhnlich um so entschiedener zeigen, je schwieriger und verwickelter die Untersuchungen sind. Inzwischen ist es doch immer von hoher Wichtigkeit, dass auch die geometrische Methode fortwährend cultivirt werde. Abgesehen davon, dass sie doch in manchen einzelnen
